%% SYNTHETIC TEST UNIT generated by writing/tools/make_outline_book.py from FINAL_OUTLINE_v2 -- not book text
\chapter{全阶因子化证明的一般方法}\label{ch:method}
本单元为测试用的合成文本（由 FINAL\_OUTLINE\_v2 自动生成），陈述内容从略。

\section{第 1 节（测试）}\label{sec:method-strategy}

\begin{definition}[区域：逐线模式指派与偏序]\label{def:method-regions}
（测试用合成陈述：内容从略。）
\uses{def:ir-pinch-surface, def:ir-scalings, thm:ir-leading-regions}
\end{definition}

\begin{lemma}[逐线交与钳夹集的交]\label{lem:method-meets}
（测试用合成陈述：内容从略。）
\uses{def:method-regions}
\end{lemma}

\begin{definition}[积分范数意义下的幂次压低]\label{def:method-power-suppressed}
（测试用合成陈述：内容从略。）
\uses{def:intro-error-types}
\end{definition}

\begin{definition}[依赖于 $\lambda$ 的邻域]\label{def:method-neighbourhoods}
（测试用合成陈述：内容从略。）
\uses{def:method-regions}
\end{definition}

\begin{lemma}[邻域覆盖]\label{lem:method-covering}
（测试用合成陈述：内容从略。）
\uses{def:method-neighbourhoods, lem:appN-beta}
\end{lemma}

\begin{definition}[共线近似子（幂等投影形式）]\label{def:method-collinear-approx}
（测试用合成陈述：内容从略。）
\uses{def:method-regions, def:prelim-lightfront, lem:prelim-boost}
\end{definition}

\section{第 2 节（测试）}\label{sec:method-regions}

\begin{definition}[软近似子（幂等投影形式）]\label{def:method-soft-approx}
（测试用合成陈述：内容从略。）
\uses{def:method-regions, def:prelim-lightfront}
\end{definition}

\begin{lemma}[软近似分母的 $i0$ 相容性与变形幅度 \tagR]\label{lem:method-i0-compat}
（$R_{\mathrm{PT}}$：在重整化微扰论中对红外正规化的部分子态严格。）
（测试用合成陈述：内容从略。）
\uses{def:method-soft-approx, def:ir-glauber}
\end{lemma}

\begin{lemma}[乘积的伸缩恒等式]\label{lem:method-telescoping}
（测试用合成陈述：内容从略。）
% uses: none
\end{lemma}

\begin{lemma}[共线近似的精度 \tagA]\label{lem:method-collinear-accuracy}
（测试用合成陈述：内容从略。）
\uses{def:method-collinear-approx, prop:ir-power-counting}
\end{lemma}
\begin{proofsketch}测试。\end{proofsketch}
\paragraph{地位与缺口}测试：本陈述的地位由\cref{lem:method-local-accuracy}、\cref{asm:appN-envelope}、\cref{asm:appN-separation}决定。

\begin{lemma}[软近似的精度 \tagA]\label{lem:method-soft-accuracy}
（测试用合成陈述：内容从略。）
\uses{def:method-soft-approx, lem:method-i0-compat, prop:ir-power-counting}
\end{lemma}
\begin{proofsketch}测试。\end{proofsketch}
\paragraph{地位与缺口}测试：本陈述的地位由\cref{asm:appN-envelope}、\cref{asm:appN-separation}决定。

\begin{lemma}[局部精度（积分范数） \tagR]\label{lem:method-local-accuracy}
（$R_{\mathrm{cond}}$：在所列编号假设之下证明完整；结论的地位不高于其中最弱的假设。）
前提：\cref{asm:appN-envelope}、\cref{asm:appN-separation}。
（测试用合成陈述：内容从略。）
\uses{def:method-power-suppressed, def:method-neighbourhoods, lem:method-telescoping, def:method-collinear-approx, def:method-soft-approx, asm:appN-envelope, asm:appN-separation, lem:appN-error-table}
\end{lemma}

\section{第 3 节（测试）}\label{sec:method-gy}

\begin{remark}[吸收性质 (Abs) 何时失效]\label{rem:method-abs-failure}
（测试用合成陈述：内容从略。）
% uses: none
\end{remark}

\begin{lemma}[Grammer–Yennie 分解]\label{lem:method-gy}
（测试用合成陈述：内容从略。）
\uses{def:method-collinear-approx, def:method-soft-approx}
\end{lemma}

\begin{proposition}[G 部分的幂次压低 \tagA]\label{prop:method-gy-suppression}
（测试用合成陈述：内容从略。）
\uses{lem:method-gy, prop:ir-power-counting, prop:ir-super-leading}
\end{proposition}
\begin{proofsketch}测试。\end{proofsketch}
\paragraph{地位与缺口}测试：本陈述的地位由\cref{asm:appN-envelope}决定。

\begin{lemma}[基本 Ward 恒等式 \tagR]\label{lem:method-ward-elementary}
（$R_{\mathrm{PT}}$：在重整化微扰论中对红外正规化的部分子态严格。）
（测试用合成陈述：内容从略。）
\uses{def:prelim-feynman-rules, prop:prelim-wilson-rules}
\end{lemma}

\begin{lemma}[阿贝尔 Ward 恒等式 \tagR]\label{lem:method-ward-abelian}
（$R_{\mathrm{PT}}$：在重整化微扰论中对红外正规化的部分子态严格。）
（测试用合成陈述：内容从略。）
\uses{lem:method-ward-elementary, lem:prelim-eikonal-identity}
\end{lemma}

\begin{proposition}[单个纵向胶子的非阿贝尔 Ward 恒等式 \tagR]\label{prop:method-ward-one}
（$R_{\mathrm{PT}}$：在重整化微扰论中对红外正规化的部分子态严格。）
（测试用合成陈述：内容从略。）
\uses{thm:prelim-slavnov-taylor}
\end{proposition}

\section{第 4 节（测试）}\label{sec:method-ward}

\begin{proposition}[树图层次的非阿贝尔恒等式与 PDF Wilson 线的种子 \tagR]\label{prop:method-ward-tree}
（$R_{\mathrm{PT}}$：在重整化微扰论中对红外正规化的部分子态严格。）
（测试用合成陈述：内容从略。）
\uses{lem:method-ward-elementary, prop:method-ward-one, prop:appF-two-gluon}
\end{proposition}

\begin{example}[$q\bar q\to g^*g^*$ 的两个纵向胶子]\label{ex:method-ward-tree}
（测试用合成陈述：内容从略。）
\uses{prop:method-ward-tree, ex:qcd-qqbar-gg}
\end{example}

\begin{lemma}[多个 K 胶子的 Ward 恒等式与 Wilson 线 \tagA]\label{lem:method-ward-multi}
（测试用合成陈述：内容从略。）
\uses{prop:method-ward-one, prop:method-ward-tree, prop:appF-induction}
\end{lemma}
\begin{proofsketch}测试。\end{proofsketch}
\paragraph{地位与缺口}测试：本陈述的地位由\cref{asm:appN-ward-product}、\cref{asm:appN-brst-nonlocal}、\cref{conj:front-brst-multilocal}、\cref{thm:method-template}、\cref{thm:appN-structure}决定。

\begin{proposition}[共线 Wilson 线与 PDF \tagA]\label{prop:method-collinear-wilson}
（测试用合成陈述：内容从略。）
\uses{lem:method-ward-multi, def:method-collinear-approx, prop:prelim-wilson-covariance, def:pdf-quark, def:pdf-gluon}
\end{proposition}
\begin{proofsketch}测试。\end{proofsketch}
\paragraph{地位与缺口}测试：本陈述的地位由\cref{lem:method-ward-multi}、\cref{asm:appN-ward-product}决定。

\begin{proposition}[软 Wilson 线 \tagA]\label{prop:method-soft-wilson}
（测试用合成陈述：内容从略。）
\uses{lem:method-ward-multi, def:method-soft-approx, lem:method-soft-accuracy}
\end{proposition}
\begin{proofsketch}测试。\end{proofsketch}
\paragraph{地位与缺口}测试：本陈述的地位由\cref{lem:method-ward-multi}、\cref{lem:method-soft-accuracy}决定。

\begin{remark}[纯规范论证的地位]\label{rem:method-pure-gauge}
（测试用合成陈述：内容从略。）
% uses: none
\end{remark}

\section{第 5 节（测试）}\label{sec:method-subtraction}

\begin{definition}[减除项]\label{def:method-subtraction}
（测试用合成陈述：内容从略。）
\uses{def:method-regions, def:method-collinear-approx, def:method-soft-approx}
\end{definition}

\begin{example}[朴素归纳为何失败：离壳单圈三角图 \tagR]\label{ex:method-naive-induction}
（$R_{\mathrm{PT}}$：在重整化微扰论中对红外正规化的部分子态严格。）
（测试用合成陈述：内容从略。）
\uses{def:method-subtraction, def:method-power-suppressed, prop:appN-triangle-exact}
\end{example}

\begin{definition}[减除条件 (CS1)–(CS3)]\label{def:method-collins-conditions}
（测试用合成陈述：内容从略。）
\uses{def:method-subtraction, def:method-power-suppressed, def:method-neighbourhoods}
\end{definition}

\begin{proposition}[减除的链展开]\label{prop:method-chains}
（测试用合成陈述：内容从略。）
\uses{def:method-subtraction}
\end{proposition}

\begin{lemma}[配对分解]\label{lem:method-pairing}
（测试用合成陈述：内容从略。）
\uses{prop:method-chains, lem:method-meets, def:appN-meet-closed}
\end{lemma}

\begin{proposition}[乘积公式]\label{prop:method-product}
（测试用合成陈述：内容从略。）
\uses{prop:method-chains, lem:method-meets, lem:method-pairing, def:method-collins-conditions, def:appN-meet-closed}
\end{proposition}

\section{第 6 节（测试）}\label{sec:method-unitarity}

\begin{lemma}[有界的亚纯函数无极点]\label{lem:method-no-pole}
（测试用合成陈述：内容从略。）
% uses: none
\end{lemma}

\begin{corollary}[去掉维数正规化 \tagR]\label{cor:method-removing-eps}
（$R_{\mathrm{cond}}$：在所列编号假设之下证明完整；结论的地位不高于其中最弱的假设。）
前提：\cref{asm:appN-contour}、\cref{thm:appN-subtraction}。
（测试用合成陈述：内容从略。）
\uses{thm:appN-subtraction, lem:method-no-pole, asm:appN-contour}
\end{corollary}

\begin{theorem}[减除定理 \tagA]\label{thm:method-subtraction}
（测试用合成陈述：内容从略。）
% errtype: E-N
\uses{def:method-subtraction, def:method-collins-conditions, def:method-power-suppressed, thm:appN-subtraction, lem:method-local-accuracy, prop:method-product, cor:method-removing-eps, asm:appN-collins, asm:appN-uv, asm:appN-contour}
\end{theorem}
\begin{proofsketch}测试。\end{proofsketch}
\paragraph{地位与缺口}测试：本陈述的地位由\cref{asm:appN-collins}决定。

\begin{example}[完整的单圈例子：离壳三角图的 Collins 减除]\label{ex:method-sudakov-subtraction}
（测试用合成陈述：内容从略。）
\uses{thm:method-subtraction, ex:method-naive-induction, prop:appN-triangle-exact, ex:sudakov-offshell, rem:ir-regions-values}
\end{example}

\begin{example}[减除语言中的单圈 DIS]\label{ex:method-dis-oneloop}
（测试用合成陈述：内容从略。）
\uses{def:method-subtraction, prop:pdf-partonic-dr, prop:pdf-pqq-operator}
\end{example}

\begin{lemma}[软因子的幺正性 \tagR]\label{lem:method-soft-unitarity}
（$R_{\mathrm{def}}$：由算符定义直接推出。）
（测试用合成陈述：内容从略。）
\uses{def:prelim-wilson, prop:prelim-wilson-covariance}
\end{lemma}

\section{第 7 节（测试）}\label{sec:method-glauber}

\begin{proposition}[末态相互作用的幺正性相消 \tagA]\label{prop:method-fsi}
（测试用合成陈述：内容从略。）
\uses{lem:method-soft-unitarity, thm:prelim-cutting-rules, lem:appO-cut-position, lem:appO-unitarity, asm:appN-cutsum, asm:appO-lcopt}
\end{proposition}
\begin{proofsketch}测试。\end{proofsketch}
\paragraph{地位与缺口}测试：本陈述的地位由\cref{asm:appN-cutsum}、\cref{asm:appO-lcopt}决定。

\begin{proposition}[两条程函线之间的单胶子交换 \tagR]\label{prop:method-glauber-oneloop}
（$R_{\mathrm{PT}}$：在重整化微扰论中对红外正规化的部分子态严格。）
（测试用合成陈述：内容从略。）
\uses{prop:prelim-wilson-rules, def:ir-glauber}
\end{proposition}

\begin{example}[单胶子交换的极点图]\label{ex:method-glauber-oneloop}
（测试用合成陈述：内容从略。）
\uses{prop:method-glauber-oneloop}
\end{example}

\begin{proposition}[Glauber 区的围道变形判据 \tagA]\label{prop:method-glauber-deformation}
（测试用合成陈述：内容从略。）
\uses{lem:method-i0-compat, prop:ir-glauber-trapping, prop:method-fsi, def:ir-glauber, prop:appO-target, thm:appO-main, asm:appO-deformation, asm:appO-mixing, asm:appO-envelope, asm:appO-rewritten-cs}
\end{proposition}
\begin{proofsketch}测试。\end{proofsketch}
\paragraph{地位与缺口}测试：本陈述的地位由\cref{asm:appO-deformation}、\cref{asm:appO-mixing}、\cref{asm:appO-envelope}、\cref{asm:appO-rewritten-cs}决定。

\begin{remark}[一个分量的形变为何不够]\label{rem:method-two-component}
（测试用合成陈述：内容从略。）
% uses: none
\end{remark}

\begin{theorem}[因子化模板定理 \tagA]\label{thm:method-template}
（测试用合成陈述：内容从略。）
非微扰地位：\cref{asm:intro-bridge}。
\uses{thm:ir-leading-regions, lem:method-ward-multi, prop:method-collinear-wilson, prop:method-soft-wilson, prop:method-glauber-deformation, lem:method-soft-unitarity, prop:method-fsi, thm:method-subtraction, prop:pdf-partonic-dr, asm:intro-bridge, asm:appN-envelope, asm:appN-separation, asm:appN-ward-product, asm:appN-cutsum, asm:appN-collins, asm:appN-values, asm:appN-uv, asm:appN-multren}
\end{theorem}
\begin{proofsketch}测试。\end{proofsketch}
\paragraph{地位与缺口}测试：本陈述的地位由\cref{thm:method-subtraction}、\cref{asm:appN-brst-nonlocal}决定。

\section{第 8 节（测试）}\label{sec:method-template}

\begin{proposition}[因子化标度与方案 \tagR]\label{prop:method-scale-scheme}
（$R_{\mathrm{PT}}$：在重整化微扰论中对红外正规化的部分子态严格。）
（测试用合成陈述：内容从略。）
（本命题只涉及部分子态。）
\uses{def:intro-factorization, thm:pdf-dglap, prop:prelim-beta, rem:pdf-schemes}
\end{proposition}
